\documentclass[12pt, letterpaper]{article}

\usepackage[utf8]{inputenc}
\usepackage[T1]{fontenc}
\usepackage[spanish]{babel}
\usepackage{geometry}
\usepackage{amsmath, amssymb, amsthm}
\usepackage{booktabs}
\usepackage{array}
\usepackage{microtype}
\usepackage{setspace}
\usepackage{parskip}
\usepackage{titlesec}
\usepackage{fancyhdr}
\usepackage{enumitem}
\usepackage{bm}
\usepackage{mathtools}
\usepackage{tcolorbox}
\usepackage{hyperref}

\tcbuselibrary{skins, breakable}

\geometry{top=2.8cm, bottom=3.0cm, left=3.2cm, right=3.2cm, headheight=14pt}
\setstretch{1.20}

\pagestyle{fancy}
\fancyhf{}
\fancyhead[C]{\small\textsc{Econometría --- Gujarati --- Ejercicio 13.11}}
\fancyfoot[C]{\small\thepage}
\renewcommand{\headrulewidth}{0.4pt}
\renewcommand{\footrulewidth}{0pt}

\titleformat{\section}
  {\large\bfseries\scshape}{\thesection.}{0.6em}{}
  [\vspace{-4pt}\rule{\linewidth}{0.4pt}]
\titleformat{\subsection}{\normalsize\bfseries}{\thesubsection.}{0.5em}{}
\titleformat{\subsubsection}{\normalsize\itshape}{}{0em}{}

\newtcolorbox{clave}{
  colback=white, colframe=black,
  boxrule=0.5pt, arc=3pt,
  left=8pt, right=8pt, top=6pt, bottom=6pt,
  breakable
}

\newcommand{\plim}{\operatorname{plim}}
\newcommand{\Cov}{\operatorname{Cov}}
\newcommand{\Var}{\operatorname{Var}}
\newcommand{\E}{\operatorname{E}}

\begin{document}

\begin{center}
  {\LARGE\bfseries Errores de Medición en la Variable Explicativa}\\[6pt]
  {\LARGE\bfseries y sus Efectos sobre los Estimadores MCO}\\[12pt]
  {\large\itshape Cambio de origen, cambio de escala y error}\\
  {\large\itshape clásico de medición (Classical Errors-in-Variables)}\\[14pt]
  \rule{0.55\linewidth}{0.6pt}\\[8pt]
  {\normalsize Ejercicio~13.11 --- \textit{Econometría}, Gujarati \& Porter, 5.a ed.}\\[4pt]
  {\small\textsc{Análisis conceptual, teórico y algebraico}}
\end{center}

\vspace{10pt}

%------------------------------------------------------------------------------
\section{Planteamiento del problema}
%------------------------------------------------------------------------------

Se considera el modelo poblacional \emph{verdadero}:
\begin{equation}
  Y_i = \beta_1 + \beta_2 X_i^* + u_i
  \label{eq:modelo_verdadero}
\end{equation}
donde $X_i^*$ es la variable explicativa \textbf{real pero no observable},
y $u_i$ es el término de perturbación que satisface las propiedades clásicas:
\[
  \E[u_i] = 0, \qquad \Cov(X_i^*, u_i) = 0, \qquad \Var(u_i) = \sigma_u^2.
\]
En la práctica, $X_i^*$ no se observa directamente. En su lugar se utiliza
un \emph{sustituto medido} $X_i$ que puede diferir de $X_i^*$ de tres formas.
El objetivo es determinar cómo cada tipo de error de medición afecta las
propiedades de los estimadores MCO de $\beta_1$ y $\beta_2$.

%------------------------------------------------------------------------------
\section{Caso a) Cambio de origen: $X_i = X_i^* + 5$}
%------------------------------------------------------------------------------

\subsection{Análisis algebraico}

Si la variable medida satisface $X_i = X_i^* + 5$, podemos despejar la
variable verdadera:
\[
  X_i^* = X_i - 5.
\]
Sustituyendo en~\eqref{eq:modelo_verdadero}:
\begin{align}
  Y_i &= \beta_1 + \beta_2\,(X_i - 5) + u_i \notag\\
      &= \beta_1 + \beta_2 X_i - 5\beta_2 + u_i \notag\\
      &= \underbrace{(\beta_1 - 5\beta_2)}_{\beta_1^*} + \beta_2 X_i + u_i.
  \label{eq:caso_a}
\end{align}
El modelo reparametrizado~\eqref{eq:caso_a} es un modelo de regresión
lineal clásico en la variable \emph{observable} $X_i$ con los supuestos
del modelo clásico plenamente satisfechos, pues el error compuesto sigue
siendo $u_i$, independiente de $X_i$.

\subsection{Efectos sobre los estimadores MCO}

\textbf{Sobre $\beta_2$ (pendiente):} El coeficiente que acompaña a $X_i$
en~\eqref{eq:caso_a} es exactamente $\beta_2$. Por tanto, el estimador
MCO $\hat{\beta}_2$ es un estimador \textbf{insesgado y consistente} del
verdadero $\beta_2$. Un cambio de origen en la variable independiente no
altera la pendiente.

\textbf{Sobre $\beta_1$ (intercepto):} El estimador MCO converge al nuevo
intercepto $\beta_1^* = \beta_1 - 5\beta_2$, \emph{no} al verdadero
$\beta_1$. Por tanto:
\[
  \E[\hat{\beta}_1] = \beta_1 - 5\beta_2 \neq \beta_1,
\]
lo que implica que $\hat{\beta}_1$ es un estimador \textbf{sesgado} de
$\beta_1$. Sin embargo, si el investigador conoce el desplazamiento (5 en
este caso), puede recuperar el verdadero intercepto mediante:
\[
  \hat{\beta}_1^{\text{corr}} = \hat{\beta}_1 + 5\hat{\beta}_2.
\]

\begin{clave}
  \textbf{Resumen del caso a):}\\
  $\hat{\beta}_2$: \textbf{insesgado y consistente} $\Rightarrow$ ningún
  problema para la pendiente.\\
  $\hat{\beta}_1$: \textbf{sesgado} en $-5\beta_2$ $\Rightarrow$ el
  intercepto original \emph{no} se recupera directamente.\\
  La relación marginal entre $Y$ y $X^*$ se preserva; solo cambia el nivel
  del intercepto.
\end{clave}

%------------------------------------------------------------------------------
\section{Caso b) Cambio de escala: $X_i = 3X_i^*$}
%------------------------------------------------------------------------------

\subsection{Análisis algebraico}

Si la variable medida es $X_i = 3X_i^*$, entonces:
\[
  X_i^* = \frac{X_i}{3}.
\]
Sustituyendo en~\eqref{eq:modelo_verdadero}:
\begin{align}
  Y_i &= \beta_1 + \beta_2\,\frac{X_i}{3} + u_i \notag\\
      &= \beta_1 + \underbrace{\frac{\beta_2}{3}}_{\beta_2^*} X_i + u_i.
  \label{eq:caso_b}
\end{align}
Nuevamente el modelo en términos de $X_i$ cumple todos los supuestos
clásicos, con $u_i$ independiente de $X_i$.

\subsection{Efectos sobre los estimadores MCO}

\textbf{Sobre $\beta_1$ (intercepto):} El intercepto no se ve afectado
por el cambio de escala. El estimador MCO satisface
$\E[\hat{\beta}_1] = \beta_1$, es decir, es \textbf{insesgado y consistente}.

\textbf{Sobre $\beta_2$ (pendiente):} El estimador MCO converge al nuevo
coeficiente de escala:
\[
  \E[\hat{\beta}_2^*] = \frac{\beta_2}{3} \neq \beta_2.
\]
El verdadero $\beta_2$ puede recuperarse mediante:
\[
  \hat{\beta}_2^{\text{corr}} = 3\,\hat{\beta}_2^*.
\]
Intuitivamente, si se mide $X^*$ en unidades tres veces más grandes, el
efecto de un incremento unitario en $X_i$ sobre $Y_i$ es la tercera parte
del efecto real sobre $X_i^*$.

\textbf{Sobre la inferencia estadística:} Los estadísticos $t$ y $F$, así
como el $R^2$ y $\bar{R}^2$, \textbf{no se ven afectados} por el cambio de
escala. La significancia estadística de los coeficientes se preserva porque
los errores estándar se escalan de la misma forma que los coeficientes,
de modo que la razón $t = \hat{\beta}/\text{ee}(\hat{\beta})$ permanece
inalterada.

\begin{clave}
  \textbf{Resumen del caso b):}\\
  $\hat{\beta}_1$: \textbf{insesgado y consistente} $\Rightarrow$ ningún
  problema para el intercepto.\\
  $\hat{\beta}_2$: \textbf{sesgado} en un factor de $1/3$ $\Rightarrow$
  subestima el efecto marginal real.\\
  Los estadísticos de prueba y el ajuste $R^2$ no cambian; sólo el valor
  numérico del coeficiente de pendiente se escala.
\end{clave}

%------------------------------------------------------------------------------
\section{Caso c) Error clásico de medición (CEV): $X_i = X_i^* + \varepsilon_i$}
%------------------------------------------------------------------------------

Este es el escenario más grave y econométricamente relevante.
Se conoce como el \textbf{Modelo Clásico de Errores en las Variables}
(\textit{Classical Errors-in-Variables}, CEV).

\subsection{Supuestos sobre el error de medición}

El término $\varepsilon_i$ es un error puramente aleatorio con las
propiedades usuales:
\[
  \E[\varepsilon_i] = 0, \qquad
  \Var(\varepsilon_i) = \sigma_\varepsilon^2, \qquad
  \Cov(X_i^*, \varepsilon_i) = 0, \qquad
  \Cov(u_i, \varepsilon_i) = 0.
\]

\subsection{Análisis algebraico}

Despejando $X_i^* = X_i - \varepsilon_i$ y sustituyendo en~\eqref{eq:modelo_verdadero}:
\begin{align}
  Y_i &= \beta_1 + \beta_2\,(X_i - \varepsilon_i) + u_i \notag\\
      &= \beta_1 + \beta_2 X_i + \underbrace{(u_i - \beta_2\,\varepsilon_i)}_{v_i}.
  \label{eq:caso_c}
\end{align}
El nuevo término de error compuesto es $v_i = u_i - \beta_2\,\varepsilon_i$.

\subsection{Violación del supuesto de exogeneidad}

El problema central surge porque $X_i = X_i^* + \varepsilon_i$ y
$v_i = u_i - \beta_2\,\varepsilon_i$ comparten el término $\varepsilon_i$.
Calculamos la covarianza entre el regresor observable y el error compuesto:
\begin{align}
  \Cov(X_i, v_i)
    &= \Cov\!\bigl(X_i^* + \varepsilon_i,\; u_i - \beta_2\,\varepsilon_i\bigr) \notag\\
    &= \Cov(X_i^*, u_i) - \beta_2\,\Cov(X_i^*, \varepsilon_i)
       + \Cov(\varepsilon_i, u_i) - \beta_2\,\Var(\varepsilon_i) \notag\\
    &= 0 - \beta_2 \cdot 0 + 0 - \beta_2\,\sigma_\varepsilon^2 \notag\\
    &= -\beta_2\,\sigma_\varepsilon^2 \neq 0.
  \label{eq:cov_no_nula}
\end{align}
Esto viola directamente el supuesto MCO de que los regresores son
ortogonales al error. El estimador MCO es, por tanto, \textbf{sesgado e
inconsistente}.

\subsection{Sesgo de atenuación y límite de probabilidad}

El límite de probabilidad del estimador MCO de la pendiente se obtiene
aplicando la ley de los grandes números:
\begin{align}
  \plim\,\hat{\beta}_2
    &= \frac{\Cov(X_i, Y_i)}{\Var(X_i)}
     = \frac{\Cov(X_i,\; \beta_1 + \beta_2 X_i + v_i)}{\Var(X_i)} \notag\\
    &= \beta_2 + \frac{\Cov(X_i, v_i)}{\Var(X_i)}
     = \beta_2 - \frac{\beta_2\,\sigma_\varepsilon^2}{\sigma_{X^*}^2 + \sigma_\varepsilon^2} \notag\\
    &= \beta_2 \left(\frac{\sigma_{X^*}^2}{\sigma_{X^*}^2 + \sigma_\varepsilon^2}\right).
  \label{eq:plim}
\end{align}
Definiendo el \emph{cociente de confiabilidad} (\textit{reliability ratio})
como:
\[
  \lambda = \frac{\sigma_{X^*}^2}{\sigma_{X^*}^2 + \sigma_\varepsilon^2}
  \in (0, 1),
\]
se obtiene:
\begin{equation}
  \plim\,\hat{\beta}_2 = \lambda\,\beta_2.
  \label{eq:atenuacion}
\end{equation}
Como $0 < \lambda < 1$, el estimador MCO converge a un valor con
\textbf{menor valor absoluto que el verdadero $\beta_2$}. Este fenómeno
se conoce como \textbf{sesgo de atenuación} (\textit{attenuation bias}):
la pendiente se «atenúa» o «encoge» hacia cero. El sesgo \emph{no desaparece}
al aumentar el tamaño muestral (es un problema de consistencia, no solo de
muestras finitas).

\subsection{Efecto sobre el intercepto}

A partir del resultado~\eqref{eq:atenuacion}, el límite de probabilidad
del estimador del intercepto es:
\begin{align}
  \plim\,\hat{\beta}_1
    &= \plim\!\bigl(\bar{Y} - \hat{\beta}_2\,\bar{X}\bigr)
     = \E[Y_i] - \lambda\,\beta_2\,\E[X_i] \notag\\
    &= (\beta_1 + \beta_2\,\mu_{X^*}) - \lambda\,\beta_2\,\mu_{X^*} \notag\\
    &= \beta_1 + \beta_2\,\mu_{X^*}\,(1 - \lambda)
     \neq \beta_1,
\end{align}
donde $\mu_{X^*} = \E[X_i^*]$. El intercepto también es \textbf{inconsistente}
como consecuencia de la inconsistencia del estimador de la pendiente. El
sesgo en $\hat{\beta}_1$ es de signo opuesto al de $\beta_2\,(1-\lambda)$.

\begin{clave}
  \textbf{Resumen del caso c):}\\
  $\hat{\beta}_2$: \textbf{sesgado e inconsistente} con
  $\plim\,\hat{\beta}_2 = \lambda\,\beta_2$, donde $\lambda \in (0,1)$.
  El estimador \emph{subestima} (en valor absoluto) el verdadero efecto.\\
  $\hat{\beta}_1$: \textbf{sesgado e inconsistente} por propagación del
  error de la pendiente.\\
  La causa raíz es la correlación no nula $\Cov(X_i, v_i) = -\beta_2\sigma_\varepsilon^2$,
  que viola el supuesto fundamental de exogeneidad del MCO.
\end{clave}

%------------------------------------------------------------------------------
\section{Síntesis comparativa}
%------------------------------------------------------------------------------

\begin{center}
\renewcommand{\arraystretch}{1.4}
\small
\begin{tabular}{p{3.6cm} p{3.5cm} p{3.5cm} p{3.0cm}}
\toprule
\textbf{Tipo de error} & \textbf{Efecto sobre $\hat{\beta}_2$} & \textbf{Efecto sobre $\hat{\beta}_1$} & \textbf{Gravedad} \\
\midrule
a)\ $X_i = X_i^* + 5$ \newline (cambio de origen)
  & Insesgado y consistente \newline ($\hat{\beta}_2 \to \beta_2$)
  & Sesgado en $-5\beta_2$ \newline (recuperable si se conoce el desplazamiento)
  & Leve \\
\midrule
b)\ $X_i = 3X_i^*$ \newline (cambio de escala)
  & Sesgado en factor $1/3$ \newline (recuperable si se conoce la escala)
  & Insesgado y consistente \newline ($\hat{\beta}_1 \to \beta_1$)
  & Leve \\
\midrule
c)\ $X_i = X_i^* + \varepsilon_i$ \newline (error aleatorio)
  & Sesgado e inconsistente; $\plim\,\hat{\beta}_2 = \lambda\beta_2 < |\beta_2|$
  & Sesgado e inconsistente
  & \textbf{Grave} \\
\bottomrule
\end{tabular}
\end{center}

\vspace{6pt}
\begin{clave}
  \textbf{Conclusión general:} Los casos a) y b) implican errores de
  medición \emph{sistemáticos y conocidos}; afectan alguno de los parámetros
  pero los efectos son recuperables. El caso c), en cambio, introduce un
  error \emph{aleatorio} que genera endogeneidad del regresor observable,
  produciendo estimadores MCO sesgados e inconsistentes para ambos parámetros.
  Los remedios habituales incluyen el uso de \textbf{Variables Instrumentales}
  (VI) o \textbf{Mínimos Cuadrados en Dos Etapas} (MC2E), que requieren
  instrumentos correlacionados con $X_i^*$ pero ortogonales a $\varepsilon_i$.
\end{clave}

\end{document}
